\chapter{Electron--phonon coupling}
\label{ch:elph}

\menu{Simulation}\menusep\menu{Electron-Phonon Coupling\dots} computes the
electron--phonon matrix elements by supercell finite differences through
GPAW's \texttt{gpaw.elph} module, and reduces them to the two numbers that make
the coupling usable downstream: the coupling constant $\lambda$ and the
electron--phonon relaxation time $\tau$ that the Drude term of the optics
module (\S\ref{sec:optics}) takes as its input.

The module is GPAW-only, and locked rather than defaulted. Every other engine's
ASE calculator exposes energies and forces; this method needs the change in the
\emph{effective potential} under an atomic displacement, which is not part of
that interface.

\section{The matrix elements}

The electron--phonon vertex is the matrix element of the potential change
produced by a phonon,
\begin{equation}
  g_{mn}^{\nu}(\mathbf{k},\mathbf{q})
    = \sqrt{\frac{\hbar}{2 M \omega_{\mathbf{q}\nu}}}\;
      \bigl\langle m,\mathbf{k}{+}\mathbf{q} \bigm|
      \boldsymbol{\epsilon}_{\mathbf{q}\nu}\cdot\nabla V
      \bigm| n,\mathbf{k} \bigr\rangle ,
  \label{eq:elph-g}
\end{equation}
where $\boldsymbol{\epsilon}_{\mathbf{q}\nu}$ is the eigenvector of mode $\nu$
at wavevector $\mathbf{q}$, $M$ the ionic mass and the square root is the
zero-point displacement amplitude. With that prefactor in place $g$ carries
units of energy; without it the bare element is in eV/\AA{} and every quantity
derived below is wrong by a mode-dependent factor. The generated script passes
\texttt{prefactor=True} for exactly this reason.

$\nabla V$ is not available analytically here: it is obtained by finite
differences. Each atom of the primitive cell is displaced by $\pm\delta$ along
$x$, $y$ and $z$ \emph{inside a supercell}, and the resulting change in the
self-consistent effective potential is cached. That is $6N+1$ supercell SCF
runs and is essentially the whole cost of the module; everything after it is
minutes.

Two constraints follow from the construction and are checked before the
expensive stage rather than after it:
\begin{enumerate}
  \item $\partial V/\partial u$ is known only out to the supercell boundary, so
    the supercell sets the \emph{range} of the interaction that can be
    represented --- and it bounds which $\mathbf{q}$ exist at all. A
    $\mathbf{q}$ that is not a reciprocal lattice vector of the supercell has
    no dynamical matrix in the cache at any price.
  \item The second delta function in \eqref{eq:a2f} is evaluated at
    $\mathbf{k}{+}\mathbf{q}$, so the electronic mesh must be a multiple of the
    phonon mesh; otherwise $\mathbf{k}{+}\mathbf{q}$ is not a sampled state and
    has no eigenvalue.
\end{enumerate}

The workflow runs in three stages: \texttt{DisplacementRunner} (the
displacements), \texttt{Supercell.calculate\_supercell\_matrix} (projection of
the potential gradients onto the LCAO basis) and
\texttt{ElectronPhononMatrix.bloch\_matrix} (rotation into the Bloch basis).
The middle stage is why the ground state must be LCAO: it projects onto
\emph{basis functions}, and a plane-wave ground state has none to project onto.

\section{The Eliashberg spectral function}

Averaging $|g|^2$ over the Fermi surface and resolving it by phonon frequency
gives the Eliashberg spectral function
\begin{equation}
  \alpha^2F(\omega) = \frac{1}{N(E_\mathrm{F})}
    \sum_{\mathbf{q}\nu}\sum_{\mathbf{k}mn}
    \bigl|g_{mn}^{\nu}(\mathbf{k},\mathbf{q})\bigr|^2
    \delta(\varepsilon_{n\mathbf{k}} - E_\mathrm{F})\,
    \delta(\varepsilon_{m\mathbf{k}+\mathbf{q}} - E_\mathrm{F})\,
    \delta(\omega - \omega_{\mathbf{q}\nu}) ,
  \label{eq:a2f}
\end{equation}
with $N(E_\mathrm{F})$ the electronic density of states at the Fermi level, in
the same per-spin per-cell normalisation as the sums. Both electron deltas sit
\emph{at} $E_\mathrm{F}$: the scattering that matters for transport connects
two states on the Fermi surface. This makes \eqref{eq:a2f} a Fermi-surface
integral, with the same slow $k$-convergence as the plasma frequency
(\S\ref{sec:optics}), and it makes the quantity meaningless for a gapped
system --- the run detects a vanishing $N(E_\mathrm{F})$ and says so rather
than returning a zero that looks like weak coupling.

The two electron delta functions are integrated by the \textbf{linear
tetrahedron method}, which interpolates the bands inside each tetrahedron
rather than requiring states to land near $E_\mathrm{F}$, and therefore has no
width to choose --- see \S\ref{sec:elph-settings} for what that replaced and
why.

The first inverse moment of \eqref{eq:a2f} is the mass-enhancement coupling
constant,
\begin{equation}
  \lambda = 2\int_0^{\infty}
    \frac{\alpha^2F(\omega)}{\omega}\,\mathrm{d}\omega ,
  \label{eq:lambda}
\end{equation}
and the logarithmic average frequency, reported alongside it because $\lambda$
alone does not say whether the coupling is to soft or stiff modes, is
\begin{equation}
  \omega_{\log} = \exp\!\left[
    \frac{2}{\lambda}\int_0^{\infty}
    \frac{\alpha^2F(\omega)}{\omega}\ln\omega \,\mathrm{d}\omega \right] .
\end{equation}

\section{From \texorpdfstring{$\lambda$}{lambda} to the Drude relaxation time}
\label{sec:elph-tau}

In the high-temperature limit the phonon-limited scattering rate is Allen's
result,
\begin{equation}
  \frac{\hbar}{\tau} = 2\pi\lambda k_\mathrm{B}T ,
  \qquad k_\mathrm{B}T \gg \hbar\omega_{\log} .
  \label{eq:allen}
\end{equation}
Two caveats are reported by the run itself, because each changes the number
rather than merely qualifying it.

\begin{description}
  \item[$\lambda$ versus $\lambda_{\mathrm{tr}}$.] Equation~\eqref{eq:allen} is
    written here with the mass-enhancement $\lambda$ of \eqref{eq:lambda}, not
    the transport coupling $\lambda_{\mathrm{tr}}$, which carries an extra
    $(1-\cos\theta_{\mathbf{k},\mathbf{k}+\mathbf{q}})$ backscattering weight
    inside the Fermi-surface average. For a free-electron-like metal such as
    aluminium or sodium the two agree to a few percent; for a transition metal
    with an anisotropic Fermi surface they need not.
  \item[The temperature range.] Below roughly $\Theta_\mathrm{D}/3$ the true
    rate crosses over to the Bloch--Gr\"uneisen $T^5$ behaviour and
    \eqref{eq:allen} \emph{overestimates} it --- so the reported $\tau$ is then
    a lower bound. The run compares the requested temperature against
    $\omega_{\log}/k_\mathrm{B}$ and warns when the limit does not apply.
\end{description}

\subsection{The factor of two}

The relaxation time is consumed by the Drude term of the optics module, and the
two conventions involved differ by a factor of two. GPAW's
\texttt{DielectricFunction} implements the free-carrier response as
\begin{equation}
  \varepsilon_{\mathrm{intra}}(\omega)
    = 1 - \frac{\omega_p^2}{(\omega + i\,\mathrm{rate})^2} ,
  \label{eq:gpaw-drude}
\end{equation}
whereas the textbook Drude function is
\begin{equation}
  \varepsilon_{\mathrm{intra}}(\omega)
    = 1 - \frac{\omega_p^2}{\omega\,(\omega + i\Gamma)} ,
  \qquad \Gamma = \frac{\hbar}{\tau} .
  \label{eq:textbook-drude}
\end{equation}
Expanding the denominator of \eqref{eq:gpaw-drude} for
$\mathrm{rate} \ll \omega$ gives
$(\omega + i\,\mathrm{rate})^2 \simeq \omega\,(\omega + 2i\,\mathrm{rate})$, so
the two agree only if
\begin{equation}
  \Gamma = 2\,\mathrm{rate}
  \qquad\Longleftrightarrow\qquad
  \boxed{\;\mathrm{rate} = \frac{\hbar}{2\tau}\;}
  \label{eq:rate-tau}
\end{equation}
GPAW's own documentation notes the same discrepancy. A $\tau$ computed here and
a \texttt{rate} quoted in a paper are therefore \emph{not} interchangeable,
which is why \texttt{epc.json} reports both $\tau$ and $\hbar/2\tau$, and why
the Optics wizard takes $\tau$ in femtoseconds and applies
\eqref{eq:rate-tau} itself.

Together with the plasma frequency --- which the $k$-point convergence module
can measure directly --- \eqref{eq:rate-tau} fixes both parameters of the Drude
model from first principles, with no fitting to the measured spectrum.

\section{Settings}
\label{sec:elph-settings}

\begin{description}
  \item[\ui{Supercell}] Repetitions the displacements are made in; the range of
    the interaction and the bound on the $\mathbf{q}$-mesh.
  \item[\ui{Displacement $\delta$}] Default 0.01\,\AA. Trades truncation error
    (anharmonicity contaminating the derivative) against numerical noise (the
    difference dominated by the SCF threshold).
  \item[\ui{LCAO basis}] \texttt{dzp} for production. Not a mode choice --- the
    projection stage requires basis functions.
  \item[\ui{Electron $k$-mesh}] Ground state and Bloch rotation. Must be a
    multiple of the $\mathbf{q}$-mesh.
  \item[\ui{Phonon $q$-mesh}] Bounded above by the supercell.
  \item[\ui{Phonon smearing}] Width of $\delta(\omega - \omega_{\mathbf{q}\nu})$,
    which bins the modes into $\alpha^2F(\omega)$ \emph{for display only}.
    $\lambda$ is summed over the modes exactly and does not change with it.
  \item[\ui{Temperature}] Where $\tau$ is evaluated. Temperature-dependent by
    physics, not by preference.
  \item[\ui{Coulomb $\mu^*$}] Morel--Anderson pseudopotential for the
    superconducting $T_\mathrm{c}$ estimate. \emph{Empirical} --- nothing here
    computes it and nothing can. Conventionally $0.10$--$0.15$ for sp metals,
    higher for transition metals.
\end{description}

\begin{calwarn}
$T_\mathrm{c}$ depends on $\mu^*$ \emph{exponentially}. For aluminium at its
literature $\lambda = 0.43$ and $\omega_{\log} = 296$~K, $\mu^* = 0.10$ gives
$1.83$~K and $0.14$ gives $0.69$~K against a measured $1.18$~K --- a factor of
$2.7$ across a range that is entirely defensible. Over the full swept range the
spread is wider still: $2.72$~K at $\mu^* = 0.08$ down to $0.36$~K at $0.16$, a
factor of $7.5$. \textbf{Quote the range, never a single value.} Every run emits
that sweep for exactly this reason.
\end{calwarn}

There is deliberately \emph{no} Fermi-surface smearing setting. The two
$\delta(\varepsilon - E_\mathrm{F})$ factors are integrated with the linear
tetrahedron method, which interpolates the bands inside each tetrahedron
rather than requiring states to land near $E_\mathrm{F}$, and so has no width
to choose; the accuracy is set by the $k$-mesh alone. Earlier versions did
expose one, and it was not cosmetic: on fcc aluminium at a $6^3$ mesh
$\lambda$ ran from $0.009$ to $31$ as that width was varied by a factor of
sixteen, with no plateau anywhere, so the reported number was whichever width
happened to be configured. Removing the parameter, rather than documenting how
to converge it, is what fixed that.

The wizard restates the cost live ($6N+1$ runs on a supercell $n\times$ the
cell) and refuses, naming the offending axis, when either commensurability
condition fails.

\section{Outputs}

The GPAW run itself stops at the raw matrix elements, writing
\texttt{elph\_raw.txt} --- a manifest naming the mesh --- beside the
eigenvalues, phonon frequencies, $\mathbf{k}+\mathbf{q}$ map and GPAW's own
\texttt{gsqklnn.npy}. Calango then integrates them. The split exists because
$|g|^2$ is (spins, $q$, $k$, modes, bands, bands) complex and reaches tens of
gigabytes on a production mesh, so the analysis is meant to run \emph{beside}
the calculation rather than after a download. Inside Calango this happens
automatically when the run finishes; on a cluster, run
\texttt{calango-elph-analyze <directory>} in place.

The analysis writes \texttt{epc.json} carrying $\alpha^2F(\omega)$ on its frequency
grid, $\lambda$, $\omega_{\log}$, $N(E_\mathrm{F})$, the relaxation time
$\tau$ at the requested temperature, the rate $\hbar/\tau$, and the
$\hbar/2\tau$ that GPAW's \texttt{rate} parameter takes. A count of imaginary
modes is included: those are a real result --- the structure is not at a local
minimum --- but $\alpha^2F$ integrates $1/\omega$ and cannot represent them, so
they are excluded and the spectral function then describes only the stable
modes.

\section{Derived properties}

Everything here comes from the same $\alpha^2F$ and costs nothing extra.

\paragraph{Superconducting $T_\mathrm{c}$.} The property $\alpha^2F$ was
invented to predict. Calango evaluates the Allen--Dynes form
\begin{equation}
  T_\mathrm{c} = f_1 f_2 \frac{\omega_{\log}}{1.2}
    \exp\!\left[\frac{-1.04(1+\lambda)}{\lambda - \mu^*(1+0.62\lambda)}\right],
\end{equation}
with the strong-coupling correction $f_1$ and the spectral-shape correction
$f_2$ (which needs the second moment $\bar\omega_2$, also reported).
McMillan's $\theta_\mathrm{D}/1.45$ form is given beside it because the older
literature quotes it, but $\omega_{\log}$ is a property of the computed
spectrum whereas $\theta_\mathrm{D}$ is a fit to a different measurement.

When $\lambda \le \mu^*(1+0.62\lambda)$ the formula has no solution: the
screened repulsion beats the phonon attraction. That is reported as ``not a
phonon-mediated superconductor at this coupling'' rather than as
$T_\mathrm{c} = 0$, which would be indistinguishable from a converged
calculation of a very small $T_\mathrm{c}$.

How well the closed form does is checked from \emph{literature} $\lambda$ and
$\omega_{\log}$ rather than from a Calango run, so what is measured is the
formula alone:

\begin{center}
\small
\begin{tabular}{@{}lllllll@{}}
\toprule
Metal & $\lambda$ & $\omega_{\log}$ & $\mu^*$ & Uncorrected & Corrected & Measured \\
\midrule
Pb & 1.55 & 56~K  & 0.10 & 6.58~K & \textbf{7.52~K} & 7.19~K \\
Al & 0.43 & 296~K & 0.12 & 1.172~K & 1.19~K & 1.18~K \\
\bottomrule
\end{tabular}
\end{center}

Lead is the case Allen and Dynes wrote $f_1$ and $f_2$ \emph{for}: at
$\lambda = 1.55$ the uncorrected form undershoots by 8\,\% and the corrections
($f_1 = 1.094$, $f_2 = 1.045$) carry it past the measurement by 5\,\%.
Aluminium, at weak coupling where $f_1 f_2 \approx 1$, lands within a per cent
either way --- but that per cent is an artefact of the $\mu^*$ chosen. The
assertion bands are correspondingly loose ($\pm 1.5$~K on Pb, $\pm 0.4$~K on
Al); the table reports what the test measures, not what it asserts.

\paragraph{Converting a bare $\mu$.} $\mu^*$ is not computed here and cannot be,
but the \emph{conversion} from a bare $\mu$ is, via Morel--Anderson
\begin{equation}
  \mu^* = \frac{\mu}{1 + \mu \ln(W/\omega_{\log})} .
\end{equation}
The retardation logarithm is computed from the \emph{occupied bandwidth}
$W = E_\mathrm{F} - \min(\varepsilon)$ --- for aluminium this came out at
11.80~eV in the reference run, against 11.66~eV for the free-electron estimate
--- and reported as \texttt{retardation\_log} beside the
\texttt{occupied\_bandwidth\_eV} it came from. The bare $\mu$ is not computed;
you supply it. A typical $\mu = 0.4$ on a typical metal lands at
$\mu^* = 0.117$, which is \emph{why} the conventional window is what it is:
$0.10$--$0.13$ for sp metals, $0.12$--$0.15$ for transition metals, where
d-electron correlation raises $\mu$.

\begin{calnote}
\textbf{The trap.} Much of the literature \emph{fits} $\mu^*$ to reproduce a
measured $T_\mathrm{c}$. That is legitimate calibration, but the resulting
$T_\mathrm{c}$ is then not a prediction. If you calibrate on a known material
and apply the value to a related one, say so.
\end{calnote}

\paragraph{Transport, and the number handed to optics.} The same Fermi-surface
sums reweighted by $1 - \cos\theta$ give the transport coupling
$\lambda_{\mathrm{tr}}$, the transport lifetime $\tau_{\mathrm{tr}}$, and the
resistivity $\rho = 1/(\varepsilon_0\omega_p^2\tau_{\mathrm{tr}})$ --- Drude
written so that every quantity in it is computed rather than assumed. Without a
plasma frequency $\rho$ is \emph{skipped, not estimated}, because
$\rho \propto 1/\omega_p^2$ and a guessed $\omega_p$ would produce a number
that looks like a measurement.

Note that $\lambda_{\mathrm{tr}}$ is \textbf{not} bounded above by $\lambda$:
since $1-\cos\theta \in [0,2]$, the only general bound is
$\lambda_{\mathrm{tr}} \le 2\lambda$. A run reporting
$\lambda_{\mathrm{tr}} > \lambda$ is telling you the $\mathbf{q}$-mesh reaches a
large fraction of $2k_\mathrm{F}$, not that something is broken.

The Drude rate handed to the optics module is
$\hbar/2\tau_{\mathrm{tr}}$ --- built on the \emph{transport} lifetime, because
a Drude term describes how a \emph{current} decays --- and it is half the
textbook damping for the reason derived in \eqref{eq:rate-tau}. A run without
band velocities falls back to the mass-enhancement $\tau$ and says so in its
warnings rather than quietly handing optics a Drude peak of the wrong width.

\paragraph{Mode-resolved coupling and linewidths.} $\lambda_{q\nu}$,
normalised so that $N_q^{-1}\sum_{q\nu}\lambda_{q\nu} = \lambda$, and the
phonon linewidth $\gamma_{q\nu} = \pi N(E_\mathrm{F})\omega_{q\nu}^2
\lambda_{q\nu}$ that follows. A single $\lambda$ cannot distinguish coupling
concentrated in one soft branch from the same total spread across the
spectrum, and $\gamma_{q\nu}$ is the one quantity in this module measurable by
inelastic neutron scattering on the same material.

\paragraph{Mass enhancement.} $m^*/m = 1 + \lambda$, and with it the
renormalisation of the linear specific-heat coefficient.

\section{Validation, and what is still wrong}

Two things are validated separately, and it matters which is which.

\paragraph{The analysis chain is validated in closed form.} It is exercised on
a contrived case with an exact answer --- a free-electron band, one phonon mode
at fixed $\omega_0$, and a constant $|g|^2$ --- so that $\lambda$ follows end to
end with nothing fitted, and a wrong normalisation, a wrong band sum or a
missing $1/N(E_\mathrm{F})$ appears as a factor rather than as a plausible
number. The same test pins $\hbar/\tau = 2\pi\lambda k_\mathrm{B}T$, the exact
factor of two in the Drude rate, the $\lambda_{\mathrm{tr}} \le 2\lambda$ bound,
and the resistivity against its closed form.

\paragraph{The shipped aluminium run is not.} The benchmark drives all three
\texttt{gpaw.elph} stages on fcc Al and asserts $0.05 < \lambda < 3.0$ --- a
window deliberately wide, because a $2\times2\times2$ supercell on a $6^3$
$k$-mesh puts only about six of 864 states within 0.1~eV of $E_\mathrm{F}$ and
cannot converge anything.

\begin{calwarn}
That run has been observed to produce $\lambda \approx 21$ against a literature
value near $0.4$ --- fifty times too large, and outside the assertion window, so
the benchmark fails on it. \textbf{The cause has not been identified.}

What is known: the analysis chain above passes its closed-form test and GPAW's
own \texttt{elph} reference test passes, so the arithmetic downstream of the
matrix elements is not the suspect. The most likely culprit is the \emph{run} ---
specifically its $2\times2\times2$ supercell, which is small enough that
aluminium comes out with \emph{imaginary acoustic branches}, and $\alpha^2F$
integrates $1/\omega$.

Two consequences for anyone using this module for a number rather than a shape.
\textbf{Validate on a material whose $\lambda$ you already know} before trusting
an absolute value from a new system; a ratio between two systems at the same
settings is far more defensible than either value alone. And \textbf{treat a
non-zero \texttt{excluded\_modes} count as a stop signal, not a footnote} --- it
says the structure is not at a local minimum in the supercell you used, so the
dynamical matrix that produced the rest of the spectrum is describing a saddle
point.
\end{calwarn}

Aluminium remains the right \emph{choice} of reference --- nearly free-electron,
$\lambda \approx 0.4$ well established, and \eqref{eq:allen} then gives
$\tau \approx 10$\,fs at room temperature, within a few femtoseconds of a Drude
fit to the measured optical conductivity. It is the particular cheap benchmark
cell that is not yet giving it.
